\documentclass[10pt,a4paper]{amsart}
\usepackage[margin=19mm]{geometry}
\usepackage{amsmath,amssymb,mathtools}
\usepackage[scr=rsfs,cal=euler]{mathalfa}
\usepackage{xfrac}
\usepackage[unicode,hidelinks,bookmarks=false]{hyperref}
\newcommand{\ie}{i.e.\ }
\newcommand{\cf}{cf.\ }
\newcommand{\resp}{respectively\ }
\newcommand{\Subsection}{\subsection}
\providecommand{\nobreakdash}{\nobreak\hbox{-}}
\setlength{\emergencystretch}{2em}
\multlinegap0pt
\usepackage{bm}
\usepackage{bbm}
\usepackage{dsfont}
\usepackage{xspace}
\usepackage{cancel}
\usepackage{mathtools}
\usepackage[shortlabels]{enumitem}
\usepackage{amsthm}
\makeatletter
\@ifundefined{theorem}{\newtheorem{theorem}{Theorem}}{}
\makeatother
\title[Relation-Theoretic Banach Contraction Principle in Topologic]{Relation-Theoretic Banach Contraction Principle in Topological Spaces with an Application}
\author{Md Hasanuzzaman, Abhishikta Das, Sumit Som}
\date{Source: arXiv:2509.10497v1 (CC BY 4.0)}
\begin{document}
\maketitle

\noindent\emph{Source and licence.} Relation-Theoretic Banach Contraction Principle in Topological Spaces with an Application. Version arXiv:2509.10497v1 (CC BY 4.0), distributed under the Creative Commons Attribution 4.0 International licence, as recorded in the arXiv licence metadata for this exact version and shown on the abstract page https://arxiv.org/abs/2509.10497v1. Full rights notice: licenses/arxiv-2509.10497-CC-BY-4.0.txt.\par\smallskip

\noindent\textit{Editorial scope.} Counted unit: the Theorem whose source label is \texttt{None} in the article (source0.tex lines 500--508, source label \texttt{None}), delivered together with the article's own complete original proof of it (source0.tex lines 509--578, one \texttt{proof} environment of 70 lines). No mathematics is written, repaired or re-derived: statement and proof are the article's own text line for line. Required context retained uncounted: main thm (lines 194--212); ode (lines 464--466). Every definition, display and label of those lines is retained unchanged. Omitted from this package and not redistributed: 1--193, 213--463, 467--499, 579--693. Nothing inside a retained excerpt is omitted or altered: every retained body is delivered exactly as the article's lines are written, so no omission receipt of any kind accompanies this package. Citations: the retained excerpts carry no citation command at all, so no bibliography entry of the article is transcribed and no reference is redistributed. Reference adaptation: none was needed and none was made; every reference inside the delivered unit resolves inside it, so no unresolved reference or citation is produced. Printed slips kept literal: no printed word, symbol, hypothesis or number of the counted statement or of its proof was altered. Layout and class: the article's own class is \texttt{amsart} with packages including amsmath, amsthm, amscd, amsfonts, amssymb, graphicx, color, epstopdf, hyperref, cite, tikz, enumerate; the wrapper is the lane's pinned \texttt{amsart} editorial wrapper at 10pt on A4, which restates the theorem-like environments the retained text needs and adds the article's own macro definitions that the retained text needs (none), with extra wrapper packages (bm, bbm, dsfont, xspace, cancel, mathtools, enumitem). The delivered numbering is the unit's own. Assets: no retained span contains a figure environment, a graphics-inclusion command, a PDF-page inclusion, a table or a TikZ picture; no image file is copied into this package, the package declares no asset dependency and no image file is redistributed with it. Licence: version arXiv:2509.10497v1 of this article is distributed under the Creative Commons Attribution 4.0 International licence, as recorded in the arXiv licence metadata for this exact version and shown on the abstract page https://arxiv.org/abs/2509.10497v1. A full rights notice is delivered at licenses/arxiv-2509.10497-CC-BY-4.0.txt. Characters: every retained excerpt is pure ASCII, so the delivered engine, XeLaTeX with its default Latin Modern text font, reports no missing character.
\begin{theorem}\label{main thm}
	 Let $\mathcal{R}$ be a binary relation on   $ \Omega $.  Suppose $ g  $  is a real-valued continuous	function on $ \Omega \times \Omega $ satisfying 
	 \begin{enumerate}[(g1)]
	 	\item $ g ( r, u ) = 0 \implies r = u $ 
	 	\item $ | g ( r, u ) | = | g ( u, r ) | $ 
	 	\item $ |g( r, u )| \leq |g(r, t ) | + | g ( t, u ) | $
	 \end{enumerate}  
     
 for all $ r, u, t  \in \Omega $ such that $ ( r, u ) \in \mathcal{R} ~ \& ~ ( t, u ) \in \mathcal{R} $   and $ S : \Omega \to \Omega $ is a mapping satisfying the followings:  
	\begin{enumerate}[(i)]
		\item $ \Omega$ is $g$-$\mathcal{R}$-complete,
		\item $\mathcal{R}$ is $ S $-closed,
		\item $ \Omega ( S;\mathcal{R})$ is non-empty,
		\item either $ S$ is ``$g-\mathcal{R}$-continuous" or ``$\mathcal{R}$ is $g$-self-closed",
		\item $ S $ is topologically $\mathcal{R}$-preserving contraction with respect to $g$.\label{iii} 
	\end{enumerate} 
    
	Then $ F (  S ) \neq \Phi $. Moreover, for each $r_0\in \Omega ( S ;\mathcal{R})$, the Picard sequence $ \{ S^n(r_0) \} $   converges to a fixed point of $S$.     
\end{theorem}

\begin{equation}\label{ode}
		D^\zeta ( f ( t ) ) = h ( t, f ( t ) ), ~~  t \in [ 0, 1 ], ~~ 1 <  \zeta \leq 2
\end{equation}

\begin{theorem}
	Consider the non-linear fractional differential equation (\ref{ode}) and suppose the function $h$ satisfies the following conditions:
	\begin{enumerate}[(i)]
		\item $ h $ is a non-decreasing function on the second variable;
		\item for each $ \mu \in [ 0, 1 ] $ and $ ( u, v ) \in \mathcal{R} $, $ h $ satisfies 
		$$ | h ( \mu, u ( \mu ) ) - h ( \mu, v ( \mu ) ) | \leq \frac{\alpha \Gamma ( \alpha + 1 ) }{4} | u ( \mu ) - v( \mu ) | ~ \text{where} ~ \alpha \in ( 0, 1 ).  $$
	\end{enumerate}
Then   the fractional differential equation (\ref{ode}) admits a solution in $C[0,1] $.
\end{theorem}

\begin{proof}
	We recall the topological space $ \Omega $,   mapping $g$ and binary relation $ \mathcal{R} $ defined above. Define a mapping $ T : \Omega \to \Omega $ by 
	\begin{equation}\label{int equn}
		T ( u( r ) ) = \dfrac{1}{\Gamma ( \zeta )} \int_{0}^{r} ( r- s)^{ \zeta - 1} h (s, u (s)) ds +   \dfrac{2 r }{\Gamma ( \zeta )} \int_{0}^{1}  \left( \int_{0}^{s} ( s - m)^{ \zeta - 1} h  ( m, u ( m ) ) dm    \right) ds~~ \qquad ~ \forall u \in \Omega.
	\end{equation}
%
Clearly the solution of (\ref{ode}) is a fixed point of $T$ in $ \Omega $. \\
Now with respect to this $T$, we verify the conditions of the Theorem \ref{main thm}. \\
%% $ \mathcal{R}$ is $T$-closed
Observe  that, for all $u, v \in \Omega $ with $ ( u, v ) \in \mathcal{R} $ and $ r \in [ 0, 1 ] $, 
\begin{align*}
	T ( u ( r ) ) = &  \dfrac{1}{\Gamma ( \zeta )} \int_{0}^{r} ( t- s)^{ \zeta - 1} h (s, u (s)) ds +   \dfrac{2 r }{\Gamma ( \zeta )} \int_{0}^{1}  \left( \int_{0}^{s} ( s - p )^{ \zeta - 1} h  ( p, u ( p ) ) dp    \right) ds \\
	\leq & \dfrac{1}{\Gamma ( \zeta )} \int_{0}^{r} ( t- s)^{ \zeta - 1} h (s, v (s)) ds +   \dfrac{2 r }{\Gamma ( \zeta )} \int_{0}^{1}  \left( \int_{0}^{s} ( s - p)^{ \zeta - 1} h  ( p, v ( p ) ) dp    \right) ds \\
    = &  T ( v ( r ) ).
\end{align*}

Hence $ ( u, v ) \in \mathcal{R} \implies   ( Tu, Tv ) \in \mathcal{R} $ and therefore $ \mathcal{R}$ is $T$-closed. 

%% $ X$ is  $g-\mathcal{R}$-complete
Next consider a $g$-Cauchy sequence $ \{ f_n\} $ in $ \Omega $. Therefore
$$  \underset{m,n \to \infty }{\lim} | g ( f_m, f_n ) | = 0  ~\qquad ~ or ~~
	  \underset{m,n \to \infty }{\lim} \underset{t \in [0, I] }{\sup} ~ |  f_m ( t ) - f_n ( t ) | = 0.  $$
	  %
      
This is the pointwise convergence of $ \{ f_n\} $ in $X$ with respect to the   metric $ d_\infty$. As the metric space $  ( \Omega,   d_\infty ) $ is complete, so $ \{ f_n\} $ converges to some $ f$ in $ ( \Omega,   d_\infty ) $. Hence $  f_n \to f $ (convergence in $ g $-sense). Moreover, if $ \{ f_n\} $ is $ \mathcal{R}$-preserving then we must have $( f_n, f)\in\mathcal{R}$ for each $n\in\mathbb{N}_0$. Thus $ \Omega $ is  $g-\mathcal{R}$-complete. 

%%% $ X (T, $\mathcal{R}$ ) \neq $\emptyset $
As $ C [ 0, 1 ] $ is non-empty, we can take   a function $ f_0 \in C [ 0, 1 ] $. Then   $ T f _0 \in C [ 0, 1 ] $. If there exist some $ r > 0 $ such that $ g ( f_0, T f _0 ) \leq r $ then $ ( f_0, T f _0 ) \in \mathcal{R} $. Otherwise choose $ f_1 \in C [ 0, 1 ] $ ($e.g.$, a function close enough to $f \in  C [ 0, 1 ] $) such that the condition satisfied. The richness of $  C [ 0, 1 ] $ ensures that such $f$ exists.  This way we can conclude $ \Omega (T; \mathcal{R} ) $ must be non-empty.    

%% $T$ is $g-\mathcal{R}$-continuous
Suppose $ \{ f_n\} $ is a $ \mathcal{R} $-preserving sequence in $ \Omega$ which converges to $ f \in \Omega $. Then 
\begin{align*} 
	& \underset{ n \to \infty }{\lim} ~ T ( f_n ) ( t ) \\
    = & \underset{ n \to \infty }{\lim} ~ \dfrac{1}{\Gamma ( \zeta )} \int_{0}^{t} ( t- s)^{ \zeta - 1} h (s, f_n (s)) ds +   \dfrac{2 t }{\Gamma ( \zeta )} \int_{0}^{1}  \left( \int_{0}^{s} ( s - m)^{ \zeta - 1} h  ( m, f_n ( m ) ) dm    \right) ds \\
	 = & \dfrac{1}{\Gamma ( \zeta )} \int_{0}^{t} ( t- s)^{ \zeta - 1} h (s, f (s)) ds +   \dfrac{2 t }{\Gamma ( \zeta )} \int_{0}^{1}  \left( \int_{0}^{s} ( s - m)^{ \zeta - 1} h  ( m, f ( m ) ) dm    \right) ds \\
     = & T ( f ) ( t )  \quad \forall t \in [ 0, 1 ]. 
\end{align*} 

Thus $T$ is $g-\mathcal{R}$-continuous.

%% $T$ is topologically $ \mathcal{R} $-preserving contraction w.r.t $g$ 
Next consider $ u, v \in \Omega $ with $ ( u, v ) \in \mathcal{R} $. Then 
$$ | g ( T u, T v ) | =  \underset{t \in [0, 1 ] }{\sup} ~ | T u ( t ) -  T v ( t ) |   $$
and  for each $ t \in [ 0, 1  ] $
\begin{align*}
	&  | T u ( t ) -  T v ( t ) |  \\
	= & ~  | \dfrac{1}{\Gamma ( \zeta )}\int_{0}^{t} ( t- s)^{ \zeta - 1} \left[  h (s, u (s) ) -  h (s, v (s) )  \right] ds + \\
		& \hskip 80 pt  \dfrac{2 t }{\Gamma ( \zeta )} \int_{0}^{1} \left(  \int_{0}^{s} ( s - p)^{ \zeta - 1} [ h  ( p, u ( p ) ) -  h  ( p, v ( p ) ) ] dp \right) ds   |   \\
		%
		\leq & ~ \dfrac{1}{\Gamma ( \zeta )} \int_{0}^{t} ( t- s)^{ \zeta - 1}  |  h (s, u (s) ) -  h (s, v (s) )  | ds + \\
		&  \hskip 80 pt  \dfrac{2 t }{\Gamma ( \zeta )} \int_{0}^{1} \left(  \int_{0}^{s} ( s - p)^{ \zeta - 1} | h  ( p, u ( p ) ) -  h  ( p, v ( p ) )  | dp   \right) ds \\
		%
		\leq & ~  \dfrac{1}{\Gamma ( \zeta )} \int_{0}^{t} ( t- s)^{ \zeta - 1}   ~ \dfrac{\alpha \Gamma ( \zeta + 1 ) }{4} | u ( s ) - v ( s ) | ds + \\
		&  \hskip 80 pt  \dfrac{2 t }{\Gamma ( \zeta )} \int_{0}^{1} \left(  \int_{0}^{s} ( s - p)^{ \zeta - 1} ~ \dfrac{\alpha  \Gamma ( \zeta + 1 ) }{4}  |   u ( p )   -    v ( p )    | dp   \right) ds \\
		%
		= & ~ \dfrac{ \alpha  \Gamma ( \zeta + 1 ) }{ 4 \Gamma ( \zeta )} \left[  \int_{0}^{t} ( t- s)^{ \zeta - 1}  | u ( s ) - v ( s ) | ds  + 2 t \int_{0}^{1} \left(  \int_{0}^{s} ( s - p )^{ \zeta - 1}    |   u ( p )   -    v ( p )    | dp   \right) ds \right]  \\
		%
		\leq & ~ \dfrac{ \alpha    \zeta   }{ 4  }  \underset{t \in [ 0, 1]}{\sup}  |   u ( t )   -    v ( t ) | \left( \int_{0}^{t} ( t- s)^{ \zeta - 1} ds + 2t  \int_{0}^{1} \left(  \int_{0}^{s} ( s - p )^{ \zeta - 1}  dp   \right) ds \right)  \\
		%
		\leq & ~ \dfrac{ \alpha  \zeta  }{ 4  } ~ | g ( u, v )|  ~ \left( \dfrac{1 + 2 t}{\zeta}   \right)  
		%
		<   ~ \alpha | g ( u, v )|. 
\end{align*}

Henceforth,
$$   | g ( T u, T v ) | \leq \alpha | g ( u, v ) |.  $$ 
%

Therefore $T$ satisfies the conditions of the  Theorem \ref{main thm} and hence $T$ has a fixed point in $ \Omega $.  Consequently   the fractional differential equation \ref{ode} admits a solution in $ C [ 0, 1 ] $. 
\end{proof}

\end{document}
